% OWNER: Arthur (Hang Seo Cho). Section draft for Milestone 5.
% Requires \usepackage{booktabs,graphicx}. Figure file: fig_loss_2x2.pdf
\section{Baseline comparison and out-of-sample evaluation}

\paragraph{Policies.}
We compare the Milestone~4 optimization model with a rule-based baseline that a
planner could apply without optimization. Using only data through June~25, the
baseline lets each store carry the $K_s=20$ products with the highest total input
demand over May~29--June~25 and orders each carried product's point forecast
$\bar d_{spt}$ (Section~2.3), with no adjustment for stock carried over from the
previous day. Days whose regional total exceeds $R_t$ are scaled down
proportionally; then, day by day, each store's quantities are scaled down until
total presale stock fits $C_{st}$ in every planning scenario. Both policies
therefore satisfy the same assortment, supply, and capacity constraints. To
separate the value of demand recovery from the value of optimization, each policy
is built twice: from raw observed sales and from stockout-adjusted demand. Both
inputs pass through the same forecast rule and scenario multipliers
$\kappa_\omega$.

\paragraph{Common evaluation.}
Every plan is scored by one simulator that applies the model's inventory rules
with the plan held fixed: zero initial stock, two-day sellable life, FIFO issuing,
sales $=\min\{\text{stock},\text{demand}\}$, and lost unmet demand. Applied to the
Milestone~4 Gurobi plan, the simulator reproduces every inventory variable and the
objective value (1{,}229.31) to within $10^{-12}$. All plans are scored with the
base-case weights $(\alpha,\beta,h,\tau)=(5,1,0.2,0.5)$ on three demand paths:
the three planning scenarios (probability-weighted), the holdout week
June~26--July~2 with stockout hours recovered from training-period references,
and the same holdout week with estimated lost demand increased by 25\%. Holdout
demand totals 1{,}106.8 normalized units (913.2 observed sales), between the low
(1{,}069.7) and typical (1{,}186.0) scenarios. No holdout information is used to
build any plan.

\paragraph{Solver note.}
For a like-for-like comparison, both optimization plans in the $2\times2$ design
are solved with the same solver and settings (HiGHS~1.16, 900-second limit,
identical formulation with 18{,}966 variables and 26{,}260 constraints). The
adjusted-input solve ends with a 2.2\% gap. The raw-input solve ends with an 11\%
gap, so its results are less precise. The official Milestone~4 Gurobi plan
(1.35\% gap) is reported alongside and gives nearly identical holdout results.

\begin{table}[h]
\centering
\small
\caption{Holdout-week performance (June 26--July 2). Quantities are normalized
units; lower loss is better. Unmet demand includes products a store does not carry.}
\label{tab:holdout}
\begin{tabular}{llrrrrrr}
\toprule
Input & Policy & Fill rate & Unmet & Expired & Carryover & Day-7 stock & Loss $J$\\
\midrule
Raw sales & Baseline     & 72.8\% & 301.4 & 10.6 & 107.5 & 33.8 & 1{,}555.8\\
          & Optimization & 73.6\% & 292.5 & 14.6 & 142.7 & 38.9 & 1{,}525.2\\
Adjusted  & Baseline     & 78.2\% & 241.6 & 19.8 & 200.8 & 58.2 & 1{,}297.2\\
          & Optimization & 78.4\% & 238.8 & 25.2 & 191.6 & 61.8 & 1{,}288.4\\
          & Optimization (Gurobi, M4) & 78.5\% & 238.3 & 26.4 & 197.9 & 63.3 & 1{,}289.1\\
\bottomrule
\end{tabular}
\end{table}

\begin{figure}[h]
\centering
\includegraphics[width=\linewidth]{fig_loss_2x2.pdf}
\caption{Weighted loss of the four policies on the three evaluation paths.}
\label{fig:loss2x2}
\end{figure}

\paragraph{Findings.}
\textbf{(1) Demand recovery matters more than optimization.} With the decision
rule fixed, switching the input from raw sales to stockout-adjusted demand lowers
holdout loss by 258.6 for the baseline and 236.8 for the optimization model (about
16\%), and raises fill rate by about 5 percentage points. Plans built on raw sales
under-order because observed sales understate demand during stockouts.

\textbf{(2) The optimization advantage shrinks out of sample.} On its own planning
scenarios, the optimization model beats the adjusted baseline by 3.3\% (Gurobi
plan: 3.7\%). On the holdout week the gap falls to 0.7\% (8.8 units of loss), and
under the +25\% recovery adjustment the two are tied (1{,}452.0 vs.\ 1{,}452.5;
the Gurobi plan is 0.3\% worse than the baseline). The result is also not uniform
across stores: on the holdout week the baseline has lower loss than the
optimization model at store~18 (370.4 vs.\ 378.0), while the optimization model is
better at stores~145 and~235. The in-sample advantage therefore partly reflects
fitting the three modeled scenarios.

\textbf{(3) Assortment drives much of the unmet demand.} For the adjusted-input
plans, 131--138 of the roughly 240 unmet units on the holdout week (55--57\%)
belong to products that a store does not carry under $K_s=20$.

\textbf{(4) The scenarios understate product-level uncertainty.} Each plan
satisfies store capacity in all three planning scenarios, which have almost no
carryover (15--33 units). On the holdout week, product-level demand departs from
the common-multiplier forecast, carryover rises to 190--201 units, and total
presale stock exceeds store capacity on 16--17 of 21 store-days, by up to about
19 units. Because all three scenarios scale every product by the same factor, they
cannot represent one product selling below forecast while another sells above it.

\subsection*{Performance under Monte Carlo demand uncertainty}

The holdout is a single week. To compare the policies over many plausible weeks,
we score every plan on the Monte Carlo demand paths of Section~\ref{sec:mc}, in which each
store-product-day varies separately (median coefficient of variation 0.40; across the 2{,}000 paths, weekly
regional demand averages 1{,}159, with a 5--95\% range of 1{,}029--1{,}298). Ten
paths are reserved for the diagnostic below, so 1{,}990 paths are used for scoring.
Because every plan is scored on the same paths, differences are paired path by
path, and we report 95\% confidence intervals.

\begin{table}[h]
\centering
\small
\caption{Mean performance over 1{,}990 Monte Carlo paths. The last column is the
paired difference from the baseline built on the same input (95\% CI).}
\label{tab:mceval}
\begin{tabular}{llrrrrr}
\toprule
Input & Policy & Fill rate & Unmet & Expired & Loss $J$ & vs.\ baseline\\
\midrule
Raw sales & Baseline     & 71.7\% & 332.0 & 3.4  & 1{,}690.1 & --\\
          & Optimization & 72.2\% & 325.2 & 11.6 & 1{,}673.9 & $-16.2$ ($-17.3$, $-15.2$)\\
Adjusted  & Baseline     & 78.4\% & 253.8 & 5.6  & 1{,}320.5 & --\\
          & Optimization & 78.4\% & 254.0 & 12.5 & 1{,}331.8 & $+11.4$ ($+10.5$, $+12.2$)\\
          & Optimization (Gurobi, M4) & 78.4\% & 253.2 & 15.0 & 1{,}330.9 & $+10.4$ ($+9.4$, $+11.3$)\\
\bottomrule
\end{tabular}
\end{table}

\begin{figure}[h]
\centering
\includegraphics[width=0.72\linewidth]{fig_milp_vs_baseline.pdf}
\caption{Loss of the optimization plan relative to the rule-based baseline
(stockout-adjusted input) on each evaluation path. The Monte Carlo bar shows the
mean paired difference with its 95\% confidence interval.}
\label{fig:milpvsbase}
\end{figure}

\textbf{(5) Demand recovery remains the dominant effect.} With the decision rule
fixed, the adjusted input lowers mean loss by 21.9\% for the baseline and 20.4\%
for the optimization model, and it lowers loss on every one of the 1{,}990 paths.

\textbf{(6) Under product-level uncertainty, the current optimization model does
not beat the baseline.} With the adjusted input, the optimization plan's mean loss
is 0.9\% higher than the baseline's, the baseline has lower loss on 70\% of
paths, and the confidence interval excludes zero. The two plans serve almost the
same demand (fill rate 78.4\% for both), but the optimization plan stocks more
heavily against the three scenarios, so it expires about twice as much stock
(12.5 vs.\ 5.6 units). Figure~\ref{fig:milpvsbase} shows the trend across
evaluation paths: the optimization advantage is 3.3\% on its own scenarios, 0.7\%
on the holdout week, zero under the +25\% test, and negative under Monte Carlo
uncertainty. With raw-sales input, the optimization plan is 0.9\% better.
Re-planning both policies on the simulated multipliers of Section~\ref{sec:mc}
(0.933, 1.020, 1.113) does not change the conclusion: the optimization plan is
still 0.8\% worse than its baseline, and both plans are worse than those built on
the original multipliers, because the lower multipliers order less.

\textbf{(7) Capacity is the binding issue.} Every adjusted-input plan exceeds store
capacity on every path, on 15--16 of 21 store-days on average (14.5 for the
baseline). As a diagnostic, we
re-solved the model with Monte Carlo paths 1--10 as ten equally weighted
scenarios (a small sample-average approximation). Constraint~(5) then requires
store capacity to hold in every one of ten paths with product-level noise. The
solver's lower bound shows what that costs: no plan that satisfies capacity in all
ten paths can have a mean loss below 1{,}454.6 on them, whereas the baseline,
which ignores capacity on paths it was not planned for, averages 1{,}247.8 on the
same ten paths. The bound is a proof from the solver, independent of the
incumbent it found. The best plan found (13.7\% gap after 900 seconds, started
from the baseline repaired to fit all ten paths) orders only 818 units, and its
mean loss on the other 1{,}990 paths is 36\% higher than the baseline's. Adding
product-level scenarios to the current model therefore makes the plan more
conservative rather than better. This suggests reconsidering how capacity is
modeled, for example by penalizing overflow or allowing excess stock to be
discarded, instead of requiring capacity to hold in every scenario.

\paragraph{Limitations.}
The holdout covers one week, so its differences below about 1\% should not be
read as meaningful; the Monte Carlo comparison gives the more reliable estimate.
The Monte Carlo paths and the holdout demand are both built from the same
stockout-recovery estimates used for planning, which favors the adjusted-input
plans; the +25\% test shows the ranking is unchanged when more lost demand is
assumed. Optimization plans were solved with time limits and end with gaps of
1.4--14\%, so part of the small differences between policies may reflect solver
precision; the raw-input and ten-path plans in particular should be re-solved
with Gurobi before final reporting.
